%%%%%%%%%%%%%%%%%%%%%%%%%%%%%%%%%%%%%%%%%%%%%%%%%%%%%%%%%%%%%%%%%%%%%%
% Lớp 10 | BPT và hệ  BPT bậc nhất hai ẩn | Bất phương trình bậc nhất hai ẩn
% Dạng 0D2?1-3: Toán thực tế  về BPT bậc nhất hai ẩn
% Số câu: 28  (sinh tự động bởi tach_cau.py; từ nay thầy sửa trực tiếp file này)
%%%%%%%%%%%%%%%%%%%%%%%%%%%%%%%%%%%%%%%%%%%%%%%%%%%%%%%%%%%%%%%%%%%%%%

%%%%%%%%%%%%%%%% TRẮC NGHIỆM 4 PHƯƠNG ÁN %%%%%%%%%%%%%%%%

% nguồn: MaHoaNH_TracNghiem_Lop10
\begin{ex}%[0D2N1-3]%[ma:df2aa2e6]
	Một cửa hàng bán lẻ bán hai loại hạt cà phê. Loại thứ nhất với  giá $140$ nghìn đồng/kg và loại thứ hai với giá $180$ nghìn đồng/kg. Cửa hàng trộn $x$ kg loại thứ nhất và $y$ kg loại thứ hai sao cho hạt cà phê đã trộn có giá không quá $170$ nghìn đồng/kg. Bất phương trình bậc nhất hai ẩn $x$, $y$ nào dưới đây mô hình hóa được yêu cầu đặt ra, với $x$, $y\in \mathbb{R}$?
	 \choice
		{$3x-y\ge 0$} 
		 {$14x+18y\ge 17$}
		 {\True $14x+18y\le 17$}
		{$3x-y\le 0$} 
	\loigiai{ 
		Vì cửa hàng trộn $x$ kg loại thứ  nhất và $y$ kg loại thứ hai sao cho hạt cà phê đã trộn có giá không quá $170$ nghìn đồng/kg nên bất phương trình tương ứng là $140x+180y\le 170\Leftrightarrow 14x+18y\le 17$.
	 }
\end{ex}

% nguồn: MaHoaNH_TracNghiem_Lop10
\begin{ex}%[0D2H1-3]%[ma:ef153027]
	Bạn Danh để dành được $900$ nghìn đồng. Trong một đợt ủng hộ trẻ em mồ côi, Danh đã lấy ra $x$ tờ tiền loại $50$ nghìn đồng, $y$ tờ tiền loại $100$ nghìn đồng  để trao tặng. Một bất phương trình mô tả điều kiện ràng buộc đối với $x,y$ là
	 \choice
		{\True  $50x+100y\le 900$}
		{$50x+100y\ge  900$}
		{$100x+50y\le 900$} 
		{$x+y=900$} 
	\loigiai{ 
		Số tiền loại $50$ nghìn đồng là $50x$ nghìn đồng.\\ 
		Số tiền  loại $100$ nghìn đồng là $100y$ nghìn đồng.\\
		Tổng số tiền bạn Danh đã lấy để trao tặng là $50x+100y$ nghìn  đồng.\\
		Mà Danh để dành  được $900$ nghìn đồng nên $50x+100y\le 900$.
	 }
\end{ex}

% nguồn: MaHoaNH_TracNghiem_Lop10
\begin{ex}%[0D2H1-3]%[ma:142fba3c]
	Bạn Lan có $15$  nghìn đồng để đi mua vở. Vở loại $A$ có giá $3\,000$ đồng một cuốn, vở loại $B$ có giá $4\,000$ đồng một cuốn. Hỏi bạn Lan có thể mua nhiều nhất bao nhiêu quyển vở sao cho bạn có cả hai loại vở?
	 \choice
		{3$ $} 
		{$5$} 
		{\True $4$} 
		 {$6$}
	 \loigiai{
	\immini{Gọi $x$, $y$ lần  lượt là số vở bạn Lan có thể mua $\left(x, y \in \mathbb{N}^{*}\right)$.\\
		Theo bài ra ta có  $3x+4y \leq 15$.\\
		Ta lấy gốc tọa  độ $O(0; 0)$ và tính $3\cdot 0+4\cdot 0-15\leq 0$.\\
		Do đó miền nghiệm của bất phương trình là nửa mặt  phẳng bờ là đường thẳng $d$ chứa gốc tọa độ $O$, kể cả đường thẳng $d$ (miền nghiệm là miền không bị gạch sọc).\\
		Vì $x$, $y \geq 1$ nên các cặp $(x; y)$ thoả mãn là $(1;1)$; $(1;2)$;  $(1;3)$; $(2;1)$; $(2;2)$; $(3;1)$.\\
		Vậy bạn Lan có thể mua được nhiều nhất $4$ quyển vở sao cho  có cả hai loại.}
		{\begin{tikzpicture}[scale=0.9,font=\footnotesize,line join=round,line cap=round,>=stealth]
		\draw[->] (-1,0)--(6,0) node[below]  {$x$};
		\draw[->] (0,-1)--(0,5)  node[left] {$y$};
		\fill  (0,0) circle (1pt) node[below left]{$O$};
		\fill[pattern=north  east lines,smooth] plot[domain=-1:6](\x,{-3/4*(\x)+15/4})-- plot[domain=6:-1](\x,{0*\x+5}) -- cycle;
		 \draw[thick,smooth,samples=200] plot[domain=-1:6] (\x,{-3/4*(\x)+15/4});
	\end{tikzpicture}
		} 
	} 
\end{ex}

% nguồn: MaHoaNH_TracNghiem_Lop10
\begin{ex}%[0D2V1-3]%[ma:525e7c24]
	Một công ty viễn thông tính phí $1$ nghìn đồng mỗi phút gọi nội mạng và $2$ nghìn đồng mỗi phút gọi ngoại mạng. Gọi $x$ và $y$ lần lượt là số phút gọi nội mạng, ngoại mạng của Bình trong một tháng và Bình muốn số tiền phải trả cho tổng đài  luôn thấp hơn $100$ nghìn đồng. Chọn khẳng định \textbf{sai}?
	 \choice
		{ Số tiền phải trả cho cuộc gọi nội mạng mỗi tháng là $x$  nghìn đồng, số tiền phải trả cho cuộc gọi ngoại mạng mỗi tháng là $2 y$ nghìn đồng. Điều kiện $x \in \mathbb{N}$, $y \in \mathbb{N}$}
		{Số tiền Bình phải trả cho tổng đài trong một tháng  là bất phương trình $x+2 y<100$ với hai ẩn $x \in \mathbb{N}$, $y \in \mathbb{N}$}
		{Cặp số $(50;  20)$ là một nghiệm của bất phương trình $x+2 y<100$}
		{\True Miền nghiệm của bất phương trình biểu diễn số tiền Bình phải trả cho tổng đài  trong một tháng là một hình vuông}
	\loigiai{ 
		\begin{itemize}
			 \item Số tiền Bình trả phí cho $x$ phút gọi nội mạng là $x$ nghìn đồng.\\
			Số tiền Bình trả phí  cho $y$ phút gọi nội mạng là $2y$ nghìn đồng.
			\item Vì trong một tháng và Bình muốn số tiền phải trả cho tổng đài luôn thấp hơn $100$ nghìn đồng nên ta có  bất phương trình $x+2y<100\,\, (\ast)$ với $x \in \mathbb{N}$, $y \in \mathbb{N}$.
			\item Ta có $50+2\cdot20=90<100$ nên  cặp số $(50; 20)$ là nghiệm của bất phương trình $x+2 y<100$.
			 \item
			 Biểu diễn miền nghiệm của $(\ast)$ trên hệ trục tọa độ. \\
			Ta thấy điểm $O(0; 0)$ thuộc miền nghiệm của $(\ast)$ do thay tọa  độ $O$ vào $(\ast)$ ta được $0<100$ (đúng).\\
			Vậy miền nghiệm của bất phương trình $(\ast)$ là nửa mặt phẳng (không kể $d$) có chứa điểm $O$  (phần không gạch chéo trên hình).
			\begin{center}
				\begin{tikzpicture}[line join=round,line cap=round,>=stealth,thick]
						 \tikzset{every node/.style={scale=0.9}}
						\begin{scope}
						\clip (-1,-1) rectangle (5,3); 
						\fill[pattern=north east lines]  (-3,3.5)--(7,3.5)--(7,-1.5)--cycle;
						\draw  (-2,3)--(6,-1);
						\end{scope}
						 \draw[->] (-1,0)--(5,0) node[below]{$x$};
						\draw[->]  (0,-1)--(0,3) node[left]{$y$};
						 \draw (0,0) node[below left]{$O$};
						\foreach  \x in {}
						\draw[thin] (\x,1pt)--(\x,-1pt) node [below]  {$\x$};
						\foreach \y in  {}
						\draw[thin] (1pt,\y)--(-1pt,\y)  node [left] {$\y$};
						\draw (4,0) node [below] {$100$} node[above]  {$B$}
						(0,2) node [below  left] {$50$} node[above right] {$A$};
					\end{tikzpicture}
			\end{center}
			Vì $x \in \mathbb{N}$, $y \in \mathbb{N}$ nên ta chỉ xét miền nghiệm bất phương trình  ứng với miền tam giác $O A B$.
		\end{itemize}
	} 
\end{ex}

% nguồn: MaHoaNH_TracNghiem_Lop10
\begin{ex}%[0D2B1-3]%[ma:72d65c94]
	Khẳng định nào sau  đây về hàm số $y=x^2$ là khẳng định đúng?
	\choice 
		{Hàm  số nghịch biến trên $\mathbb{R}$}
		{Hàm số đồng biến  trên $\mathbb{R}$}
		{Hàm số nghịch biến trên  $[0;+\infty)$}
		{\True Hàm  số đồng biến trên $[0;+\infty)$}
	\loigiai{ 
		Hàm số $y=x^2$ có hệ số $a=1>0$ nên hàm đồng biến trên $\left(-\dfrac{b}{2a};+\infty \right) $ hay $[0;+\infty)$ (vì hàm số chỉ bằng $0$  tại một điểm là $x=0$).
	 }
\end{ex}

% nguồn: MaHoaNH_TracNghiem_Lop10
\begin{ex}%[0D2K1-3]%[ma:28734e74]
	Cho tập hợp $A=\{0;1;2;3;5\}$. Từ tập $A$ có thể lập được bao nhiêu số chia hết cho $5$ gồm $4$  chữ số khác nhau?
	 \choice
		 {$54$}
		 {$72$}
		{$69$} 
		{\True  $42$}
	\loigiai{ 
		Xét số $\overline{abcd} \;\vdots\; 5$, với  $a,b,c,d$ khác nhau lấy từ $A=\{0;1;2;3;5\}$. Khi đó $d=0$ hoặc $d=5$.\\
		\begin{itemize}
			\item Nếu $d=0$ ta xét số $\overline{abc0}$, với $a,b,c \in  \{1;2;3;5\}$.\\
			Số cách chọn $a,b,c$ khác nhau (theo  thứ tự đó) là $4,3,2$.\\
			Trường hợp này có  $4\cdot 3\cdot 2=24$ số.
			\item Nếu  $d = 5$ ta xét số $\overline{abc5}$, với $a,b,c \in \{0;1;2;3\}$, $a \neq 0$.\\
			Số cách chọn $a,b,c$ khác nhau (theo thứ tự  đó) là $3,3,2$.\\
			Trường hợp này có  $3\cdot 3\cdot 2=18$ số.
		\end{itemize}
		Vậy có $24+18=42$  số thỏa mãn yêu cầu của bài toán.
	} 
\end{ex}

% nguồn: MaHoaNH_TracNghiem_Lop10
\begin{ex}%[0D2B1-3]%[ma:74e1f31f]
	\immini[thm]{Cho hàm số $y=f(x)$ có tập xác định là $[-3;3]$ và đồ thị của nó được biểu diễn bởi hình bên. Khẳng định nào  sau đây đúng?
	\choice 
		{\True Hàm  số đồng biến trên khoảng $(-3;-1)$}
		{Hàm số đồng biến  trên khoảng $(-3;3)$}
		 {Hàm số đồng biến trên khoảng $(-3;0)$}
		{Hàm số nghịch biến trên  khoảng $(-1;2)$}}{
		\begin{tikzpicture}[>=stealth,x=1.0 cm,y=1.0 cm,scale=0.75]
				\draw[->] (-3.5,0)--(4,0)  node[below left] {$x$};
				\draw[->] (0,-2)--(0,4.5) node[below right] {$y$}; 
				\foreach \x in {,-3,-2,-1,1,2,3,,} 
				\draw[shift={(\x,0)},color=black]  (0pt,2pt) -- (0pt,-2pt) node[below] {\footnotesize $\x$};
				 \foreach \y in {,-1,1,2,3,4,}
				\draw[shift={(0,\y)},color=black]  (2pt,0pt) -- (-2pt,0pt) node[left] {\footnotesize $\y$};
				\draw[color=black]  (0pt,-10pt) node[right] {\footnotesize $0$};
				 \clip (-4,-2.5) rectangle (4,4.5);
				 \draw[dashed] (0,-1)--(-3,-1)--(-3,0);
				\draw[dashed]  (-1,0)--(-1,1)--(0,1);
				\draw[dashed] (0,4)--(3,4)--(3,0); 
				\draw[magenta,thick] (-3,-1)--(-1,1)--(0,1)--(3,4); 
			\end{tikzpicture}
		 }
	\loigiai{ 
		Trên khoảng $(-3;-1)$ và $(1;3)$, đồ thị hàm số đi lên, do đó hàm số đồng biến trên $(-3;-1)$ và $(1;3)$. 
	 }
\end{ex}


%%%%%%%%%%%%%%%% ĐÚNG/SAI %%%%%%%%%%%%%%%%

% nguồn: MaHoaNH_DungSai_Lop10
\begin{ex}%[0D2H1-3]%[ma:38e0aeca]
	Một cửa hàng dành tối đa $10$ triệu để mua $x$ tạ gạo  và $y$ tạ mì. Biết mỗi tạ gạo mua giá $1{,}5$ triệu, mỗi tạ mì mua giá $1{,}2$ triệu. Khi đó
	\choiceTF 
		{\True Số tiền để mua $x$ tạ gạo và $y$ tạ mì  là $1{,}5 x+1{,}2 y$ triệu đồng}
		{Bất phương trình biểu thị mối liên hệ giữa $x$ và $y$  là $1{,}5 x+1{,}2 y \geq 10$}
		{Cửa hàng có thể mua được $5$ tạ gạo và $5$  tạ mì}
		{\True Bất phương trình biểu thị số tiền cửa hàng mua gạo và mì là $1{,}5  x+1{,}2 y \leq 10$ với $x$, $y\ge 0$}
	 \loigiai{
		\begin{itemchoice}
			\itemch  Số tiền cửa hàng mua $x$ tạ gạo là $1{,}5x$ triệu đồng.\\
			Số tiền cửa hàng mua $y$ tạ gạo là  $1{,}2y$ triệu đồng.\\
			Số tiền để mua $x$ tạ gạo và $y$ tạ mì là $1{,}5 x+1{,}2  y$ triệu đồng.
			\itemch Vì cửa hàng dành tối đa $10$ triệu để mua $x$ tạ gạo và $y$ tạ mì nên bất phương trình biểu thị mối liên hệ giữa $x$ và  $y$ là $1{,}5 x+1{,}2 y \leq 10$.
			\itemch Vì số tiền cửa hàng có thể mua $5$ tạ gạo và $5$ tạ  mì là $1{,}5 \cdot 5+1{,}2\cdot5=13{,}5>10$.\\
			Do đó cửa hàng không thể mua được $5$ tạ gạo và $5$ tạ mì. 
			\itemch Vì cửa hàng dành tối đa $10$ triệu  để mua $x$ tạ gạo và $y$ tạ mì nên bất phương trình biểu thị số tiền cửa hàng mua gạo và mì là $1{,}5 x+1{,}2 y \leq 10$ với $x$, $y\ge 0$.
		\end{itemchoice}
	 }
\end{ex}

% nguồn: MaHoaNH_DungSai_Lop10
\begin{ex}%[0D2H1-3]%[ma:5b02a6fd]
	Một công ty viễn thông tính phí $1$ nghìn đồng mỗi phút gọi nội mạng và 2nghìn đồng mỗi phút gọi ngoại mạng. Gọi $x$ và $y$ lần lượt là số phút gọi nội mạng, ngoại mạng của Bình trong một tháng. Bình muốn số tiền phải trả cho tồng đài luôn thấp hơn $100$ nghìn đồng. Khi  đó
	 \choiceTF
		{\True Số tiền phải trả cho cuộc gọi nội mạng mỗi tháng là $x$ (nghìn đồng), sô tiển phải trả cho cuộc gọi ngoại mạng mỗi tháng là $2y$ (nghìn đồng) với điều kiện: $x \in  \mathbb{N}$, $y \in \mathbb{N}$}
		{\True Bất phương trình bậc  nhất gồm hai ẩn số $x$, $y$ đã cho là $x+2y < 100$}
		{\True $x=50$, $y=20$ nghiệm của bất phương trình bậc nhất gồm hai ẩn số  $x$, $y$ đã cho}
		{Miền nghiệm của bất phương trình bậc nhất gồm hai ẩn số $x$, $y$  đã cho là một hình vuông}
	\loigiai{ 
		\begin{itemchoice}
			\itemch Số tiền phải trả cho cuộc gọi nội mạng mỗi tháng là $x$ (nghìn đồng), số tiền phải trả cho cuộc gọi ngoại mạng mỗi tháng là $2y$ (nghìn đồng) với điều kiện $x \in \mathbb{N}$, $y \in \mathbb{N}$. 
			\itemch  Ta có bất phương trình $x+2y < 100$ (*).
			\itemch Xét $x=50, y=20$ thay vào $(*)$ ta  được $50+2\cdot 20< 100$ (đúng) suy ra ($50; 20$) là một nghiệm của (*).
			\itemch  \immini{Biểu diễn miền nghiệm của (*) trên mặt phẳng tọa độ: Vẽ đường thẳng $x+2y=100$.\\
			Ta thấy điểm  $O(0; 0)$ thuộc miền nghiệm của (*) do thay tọa độ $O$ vào $(*): 0< 100$ (đúng).\\
			Vậy miền nghiệm của bất phương trình $(*)\colon x+2y < 100$ là nửa mặt phẳng (không kể d) có chứa điểm $O$ (phần không gạch chéo trên  hình).\\
			Trong thực tế, vì $x \in \mathbb{N}$, $y \in \mathbb{N}$ nên ta chỉ xét  miền nghiệm bất phương trình ứng với miền tam giác $OAB$ mà thôi.}
			{\begin{tikzpicture}[scale=0.9,font=\footnotesize,line join=round,line cap=round,>=stealth]
		\draw[->] (-2,0)--(4,0)  node[below] {$x$};
		 \draw[->] (0,-1)--(0,3) node[left] {$y$};
		\fill (0,0) circle (1pt)  node[below left]{$O$};
		\fill 
			(0,1) circle  (1pt) node[left]{$50$}
			(0,1) circle  (1pt) node[right]{$A$}
			(2,0) circle (1pt) node[below]{$100$} 
			(2,0)  circle (1pt) node[above]{$B$};
		\fill[pattern=north east lines,smooth] plot[domain=-2:4](\x,{-1/2*(\x)+1})-- plot[domain=4:-2](\x,{0*\x+3}) -- cycle; 
		\draw[thick,smooth,samples=200]  plot[domain=-2:4] (\x,{-1/2*(\x)+1});
	\end{tikzpicture}}
		\end{itemchoice}
	} 
\end{ex}

% nguồn: MaHoaNH_DungSai_Lop10
\begin{ex}%[0D2H1-3]%[ma:ac43fdbe]
	Nhân ngày Quốc tế Thiếu nhi ngày $01$ tháng $06$, một rạp chiếu phim phục vụ các khán giả một bộ phim hoạt  hình. Vé được bán ra có hai loại:
	\begin{itemize}
		\item Loại $1$ (dành cho trẻ từ $6$ từ $13$ tuổi): $50\,000$  đồng/vé.
		\item Loại $2$ (dành cho người trên $13$  tuổi): $100\,000 $ đồng/vé.
	\end{itemize}
	Người ta tính toán rằng, để không phải bù lỗ thì  số tiền vé thu được ở rạp chiếu phim này phải đạt tối thiểu $20$ triệu đồng. Gọi $x$ là số lượng vé loại $1$ bán được $(x \in \mathbb{N})$ và $y$ là số lượng vé loại $2$ bán được $(y \in \mathbb{N})$.
	 \choiceTF
		{\True Người ta sẽ phải bù lỗ trong trường hợp số tiền bán  vé thoả mãn bất phương trình $x+2y < 400$}
		 {Nếu bán được $250$ vé loại $1$ và $150$ vé loại $2$ thì rạp chiếu phim có lãi}
		{Nếu bán được $200$ vé loại $1$ và $100$ vé loại $2$ thì  rạp chiếu phim hoà vốn}
		{Nếu bán được $50$ vé  loại $1$ và $100$ vé loại $2$ thì rạp chiếu phim phải bù lỗ}
	\loigiai{Gọi $x$ là số lượng vé loại 1bán được $(x \in \mathbb{N})$ và $y$ là số lượng vé loại 2bán được $(y \in \mathbb{N})$ thì số tiền bán vé thu được là $50x+100y$ (nghìn đồng).\\ 
		Người ta sẽ phải bù lỗ trong trường hợp số tiền bán  vé nhỏ hơn $2$0triệu đồng, tức là $50x+100y < 20\,000$ hay $x+2y < 400$.\\
		Như vậy, việc giải quyết bài toán mở đầu dẫn đến việc đi tìm miền nghiệm của bất phương trình $x+2y <  400$.\\
		Miền nghiệm của  bất phương trình bậc nhất hai ẩn này được xác định như sau:\\
		Vẽ đường thẳng $d: x+2y=400$. Ta lấy gốc tọa độ $O(0; 0)$  và tính $0+2\cdot 0=0< 400$.\\
		Do đó, miền nghiệm của bất phương trình  là nửa mặt phẳng bờ $d$ chứa gốc tọa độ không kể đường thẳng $d$.
		\begin{center}
			\begin{tikzpicture}[scale=0.9,font=\footnotesize,line join=round,line cap=round,>=stealth]
					\draw[->] (-2,0)--(6,0) node[below] {$x$}; 
					\draw[->] (0,-1)--(0,3) node[left]  {$y$};
					\fill (0,0)  circle (1pt) node[below left]{$O$};

					 \fill
						 (0,1) circle (1pt) node[left]{$100$}
						(0,2)  circle (1pt) node[left]{$200$}
						(2,0) circle  (1pt) node[below]{$200$}
						(4,0) circle (1pt)  node[below]{$400$};
					\draw[dashed] (2,0)--(2,1)--(0,1); 
					\fill[pattern=north east lines,smooth]  plot[domain=-2:6](\x,{-1/2*(\x)+2})-- plot[domain=6:-2](\x,{0*\x+3}) -- cycle;
					 \draw[thick,smooth,samples=200] plot[domain=-2:6] (\x,{-1/2*(\x)+2});
				\end{tikzpicture}
		\end{center}
		\begin{itemize}
			\item Vậy nếu bán được số vé loại $1$ là $x$ và số vé loại $2$ là $y$ mà điểm $(x; y)$ nằm trong miền tam giác $OAB$ không hề cạnh $AB$ thì rạp chiếu phim sẽ phải bù  lỗ.
			 \item Nếu điểm $(x; y)$ nằm trên đoạn thẳng $AB$ thì rạp chiếu phim hoà vốn.
			\item Nếu bán được $150$ vé loại  1và $150$ vé loại $2$ thì rạp chiếu phim có lãi.
			\item Nếu bán được $200$ vé loại 1và $100$ vé loại  $2$ thì rạp chiếu phim hoà vốn.
			 \item Nếu bán được $100$ vé loại 1và $100$ vé loại $2$ thì rạp chiếu phim phải bù lỗ.
		\end{itemize}
		 Khi đó
		\begin{itemchoice}
			\itemch Người ta sẽ phải bù lỗ trong trường hợp số tiền bán vé thoả mãn bất phương trình $x+2y < 400$. 
			 \itemch Nếu bán được $200$ vé loại $1$ và $100$ vé loại $2$ thì rạp chiếu phim hoà vốn.
			\itemch Nếu bán được $200$ vé loại $1$ và  $100$ vé loại $2$ thì rạp chiếu phim hoà vốn.
			\itemch Nếu bán được $100$ vé  loại $1$ và $100$ vé loại $2$ thì rạp chiếu phim phải bù lỗ.
		\end{itemchoice}
	} 
\end{ex}


%%%%%%%%%%%%%%%% TRẢ LỜI NGẮN %%%%%%%%%%%%%%%%

% nguồn: MaHoaNH_TraLoiNgan_Lop10
\begin{bt}%[0D2H1-3]%[ma:f7d6daa7]
	Bạn Lan mang $150\,000$ đồng đi nhà sách để mua một số quyển tập và bút. Biết rằng giá một quyển tập là $8\,000$ đồng và giá của một cây bút là $6\,000$ đồng. Bạn Lan có thể mua được tối đa bao nhiêu quyển tập  nếu bạn đã mua $10$ cây bút.
	 \shortans{11}
	\loigiai{Bất phương trình biểu diễn số tập và bút  có thể mua được phụ thuộc vào số tiền mang theo là $8\,000x+6\,000y \leq 150\,000$.\\
		Bạn Lan có thể mua được tối đa số quyển tập nếu bạn đã mua $10$ cây bút là \[8\,000x+6\,000\cdot 10\leq 150\,000\Leftrightarrow x  \leq 11{,}25.\]
		Vì $x$ nguyên dương nên số quyển tập tối đa bạn Lan  mua được là $11$ quyển.
	} 
\end{bt}

% nguồn: MaHoaNH_TraLoiNgan_Lop10
\begin{bt}%[0D2H1-3]%[ma:5b2f7016]
	Một gian hàng trưng bày bàn và ghế rộng $60$ m$^2$. Diện tích để kê một chiếc ghế là $0{,}5$ m$^2$, một chiếc bàn là $1{,}2$ m$^2$. Gọi $x$ là số chiếc ghế, $y$ là số chiếc bàn được kê. Biết diện tích mặt sàn dành cho lưu thông tối thiểu là $12$ m$^2$. Giả sử gian hàng đã kê $10$ chiếc  bàn thì phần diện tích cho phép còn lại có thể kê được nhiều nhất bao nhiêu chiếc ghế?\\
	\shortans{72} 
	 \loigiai{
		Diện tích cho phép trưng bày bàn và ghế là $60-12=48$ (m$^2$).\\ 
		Ta  có bất phương trình $0{,}5x+1{,}2y\le 48$.\\
		Do đã kê $10$ chiếc bàn nên  $0{,}5x+1{,}2\cdot 10\le 48\Leftrightarrow x\le 72$.
	} 
\end{bt}

% nguồn: MaHoaNH_TraLoiNgan_Lop10
\begin{bt}%[0D2H1-3]%[ma:ce5588ed]
	Một hộ nông dân định trồng đậu và cà trên diện tích $800$ m$^2$. Nếu trồng đậu trên diện tích $100$ m$^2$ thì cần $20$ công làm và thu được $3\,000\,000$  đồng. Nếu trồng cà thì trên diện tích $100$ m$^2$ cần 3$ 0$ công làm và thu được $4\,000\,000$ đồng. Hỏi cần trồng mỗi loại cây trên diện tích là bao nhiêu để thu được nhiều tiền nhất khi tổng số công làm không quá $180$ công.
	 \shortans{200}
	 \loigiai{
		Giả sử diện tích trồng đậu là $x$,  suy ra diện tích trồng cà là $8-x$.\\
		Ta có thu nhập thu  được là $S(x)=[3x+4(8-x)] \cdot 10\,000=10\,000(-x+32)$ đồng.\\
		Tổng số công  là $20x+30(8-x)=-10x+240$.\\
		Theo giả thiết có $-10x+240\leq 180\Leftrightarrow x \geq 6$.\\ 
		Mà hàm số  $S(x)$ là hàm nghịch biến trên $\mathbb{R}$ nên $S(x)$ đạt giá trị lớn nhất khi $x=6$.\\
		Do đó trồng $600$ m$^2$ đậu và $200$  m$^2$ cà.
	 }
\end{bt}

% nguồn: MaHoaNH_TraLoiNgan_Lop10
\begin{bt}%[0D2H1-3]%[ma:35cbe374]
	Anh An là nhân viên bán hàng tại siêu thị điện máy. Anh An kiếm được một khoản hoa hồng $600$ nghìn đồng cho mỗi máy giặt và $1{,}3$ triệu đồng  cho mỗi tủ lạnh mà anh ấy bán được. Hỏi để nhận được từ $10$ triệu đồng trở lên tiền hoa hồng thì anh An cần bán bao nhiêu máy giặt và tủ lạnh?
	\shortans{} 
	\loigiai{ 
		Gọi $x$ và $y$ lần lượt là  số máy giặt và số tủ lạnh anh An bán được. Khi đó số tiền hoa hổng mà anh An nhận được là $0{,}6x+1{,}3y$ (triệu đồng).\\
		Theo để bài, ta  có: $0{,}6x+1{,}3y \geq 10$\\
		Tiếp  theo ta xác định miền nghiệm của bất phương trình $0{,}6x+1{,}3y \geq 10$ như sau:
		\begin{itemize}
			\item Bước $1$: Vẽ đường  thẳng $d: 0{,}6x+1{,}3y=10$ trên mặt phẳng toạ độ $Oxy$.
			\item Bước $2$: Lấy  điểm $O(0; 0)$ không thuộc $d$ và thay vào biểu thức $0{,}6x+1{,}3y$ ta được $0{,}6.0+1{,}3.0=0< 10$.\\
			Do đó, miền nghiệm của bất phương trình là nửa mặt phẳng bờ $d$ không chứa gốc  tọa độ (miển không bị tô màu).
			\begin{center}
				\begin{tikzpicture}[scale=0.9,font=\footnotesize,line join=round,line cap=round,>=stealth]
						\draw[->] (-1,0)--(6,0) node[below] {$x$}; 
						\draw[->] (0,-1)--(0,3) node[left] {$y$}; 
						\fill (0,0) circle (1pt)  node[below left]{$O$};
						\fill[pattern=north east lines,smooth]  plot[domain=-1:6](\x,{-1/2*(\x)+5/2})-- plot[domain=6:-1](\x,{0*\x-1}) -- cycle;
						\draw[thick,smooth,samples=200]  plot[domain=-1:6] (\x,{-1/2*(\x)+5/2});
					\end{tikzpicture}
			\end{center}
		\end{itemize}
		Vậy nếu anh An bán được số máy giặt là $x$ ($x \in \mathbb{N}$) và số tủ lạnh là $y$  ($y \in \mathbb{N}$) sao cho điểm $(x; y)$ nằm trong nửa mặt phẳng bờ $d$ không chứa gốc tọa độ thì anh An nhận được từ $10$ triệu đồng trở lên tiền hoa hồng.
	 }
\end{bt}

% nguồn: MaHoaNH_TraLoiNgan_Lop10
\begin{bt}%[0D2V1-3]%[ma:02579a34]
	Bạn Nam tiết kiệm được $600$ nghìn đồng. Trong đợt ủng hộ các bạn học sinh đồng bào miền Trung bị lũ lụt vừa qua, bạn Nam đã ủng hộ $x$ tờ tiền loại $30$ nghìn đồng, $y$ tờ tiền loại $20$ nghìn đồng. Khi đó bất phương trình biểu diễn tổng số tiền  mà bạn Nam đã ủng hộ có dạng $ax+by\le c$. Tính giá trị của biểu thức $P=c-2a-b$.
	 \shortans{490}
	 \loigiai{
		Tổng số tiền bạn Nam đã ủng  hộ là $30x+20y$.\\
		Vì bạn Nam chỉ có tất cả $600$ nghìn đồng nên tổng số tiền bạn Nam đã ủng hộ không thể vượt quá $600$ nghìn đồng. Vậy  ta có bất phương trình $30x+ 20y\le 600$.\\
		 Do đó $a=30$, $b=20$, $c=600 \Rightarrow P=600-2\cdot30-20=490$.
	} 
\end{bt}

% nguồn: MaHoaNH_TraLoiNgan_Lop10
\begin{bt}%[0D2V1-3]%[ma:14b86d28]
	Bạn Nam tiết kiệm được $450$ nghìn đồng. Trong đợt ủng hộ các  bạn học sinh đồng bào miền Trung bị lũ lụt vừa qua, bạn Nam đã ủng hộ $x$ tờ tiền loại $20$ nghìn đồng, $y$ tờ tiền loại $10$ nghìn đồng. Khi đó bất phương trình biểu diễn tổng số tiền mà bạn Nam đã ủng hộ có dạng $mx+y\le n$. Tính giá trị của biểu thức $P=m+n$.
	\shortans{47} 
	 \loigiai{
		Tổng số tiền bạn Nam đã ủng hộ là  $20x + 10y$.\\
		Vì bạn Nam chỉ có tất cả $450$ nghìn đồng nên tổng số tiền bạn Nam đã ủng hộ không  thể vượt quá $450$ nghìn đồng.\\
		Vậy ta có bất  phương trình $20 x+10 y \leq 450\Leftrightarrow 2x+y\le 45$.\\
		 Do đó $m=2$, $n=45$. Suy ra $P=m+n=47$.
	} 
\end{bt}

% nguồn: MaHoaNH_TraLoiNgan_Lop10
\begin{bt}%[0D2V1-3]%[ma:763de39a]
	Một đội sản xuất cần $3$ giờ để làm xong sản phẩm loại I và $2$ giờ để làm xong sản phẩm loại II. Biết thời gian tối đa cho việc sản xuất hai sản phẩm trên là $18$ giờ. Gọi $x$, $y$ lần lượt là số sản phẩm loại I, loại II mà đội làm được trong thời gian cho phép. Khi đó bất phương trình biểu diễn tổng thời gian làm xong hai  loại sản phẩm trên có dạng $3x+by+c\le 0$. Tính giá trị của biểu thức $P=b+c$.
	\shortans{-16} 
	\loigiai{ 
		Tổng thời gian làm xong hai loại  sản phẩm là $3x+ 2y$ giờ.\\
		Vì thời gian tối đa cho việc sản xuất hai sản phẩm trên là $18$ giờ, 
		Vậy ta có bất phương trình $3x+2y \le 18\Leftrightarrow  3x+2y -18\le 0$.\\
		Do đó $b=2$, $c= -18$. Suy  ra $P=b+c=-16$.
	 }
\end{bt}

% nguồn: MaHoaNH_TraLoiNgan_Lop10
\begin{bt}%[0D2V1-3]%[ma:4363468e]
	Một học sinh đem $300\,000$ đồng đi nhà sách để mua một số sách và bút. Biết rằng 
	giá một quyển sách là $20\,000$ đồng và  giá của một cây bút là $12\,000$ đồng. Học sinh đó có thể
	mua được tối đa bao nhiêu quyển sách nếu bạn  đã mua $10$ cây bút.
	 \shortans{9}
	\loigiai{ 
		Gọi $x$, $y$ lần lượt  là số sách và số bút mà học sinh mua $(x,y\in\mathbb{N})$.\\
		Từ giả thiết, ta có  $20\,000x+12\,000y\leq300\,000$.\\
		Học sinh  có thể mua được tối đa số quyển sách nếu bạn đã mua $10$ cây bút là $20\,000x+12\,000\cdot10\leq300\,000\Leftrightarrow x\leq9$.\\
		Vì $x$ nguyên dương nên số  quyển sách tối đa học sinh mua được là 9 quyển.
	 }
\end{bt}


%%%%%%%%%%%%%%%% TỰ LUẬN %%%%%%%%%%%%%%%%

% nguồn: MaHoaNH_TuLuan_Lop10
\begin{bt}%[0D2H1-3]%[ma:be3ff40a]
	Bạn Minh để dành được $900$ nghìn đồng. Trong một đợt ủng hộ trẻ em mồ côi, Minh đã lấy ra $x$ tờ tiền loại $50$ nghìn đồng, $y$ tờ tiền  loại $100$ nghìn đồng để trao tặng. Một bất phương trình mô tả điều kiện ràng buộc đối với $x$, $y$ là $ax+by\le c$. Tính $c-a-b$.
	 \loigiai{
		Tổng số tiền bạn Minh đã lấy để  trao tặng: $50x+100y$ (nghìn đồng).\\
		Mà Minh để dành  được $900$ nghìn đồng nên $50x+100y\le 900$.\\
		 Vậy $a=50$, $b=100$, $c=900$ nên $c-a-b=750$
	 }
\end{bt}

% nguồn: MaHoaNH_TuLuan_Lop10
\begin{bt}%[0D2H1-3]%[ma:cc34fd31]
	Một gian hàng trưng bày bàn và ghế rộng $60$ m$^2$. Diện tích để kê một chiếc ghế là $0{,}5$ m$^2$ và một chiếc bàn là $1{,}2$ m$^2$. Gọi $x$  là số chiếc ghế, $y$ là số chiếc bàn được kê.
	\begin{enumerate}
		\item Viết bất phương trình bậc nhất hai ẩn $x$, $y$ cho phần mặt sàn để kê bàn và ghế, biết diện tích mặt sàn dành cho lưu thông tối  thiểu là $12$ m$^2$.
		\item Chỉ ra ba nghiệm của  bất phương trình trên.
	\end{enumerate}
	\loigiai{\begin{enumerate}
		\item Diện tích để kê $x$  chiếc ghế, $y$ chiếc bàn là $0{,}5x+1{,}2y$ (m$^2$).\\
		Diện tích tối đa để kê bàn và  ghế là $60-12=48$ (m$^2$).\\
		Do đó ta có bất phương trình $0{,}5x+1{,}2y \le 48$. 
		\item Ba nghiệm có thể  chỉ ra được của bất phương trình trên là $(20;30)$, $(30;20)$, $(50;15)$.
\end{enumerate}}
\end{bt}

% nguồn: MaHoaNH_TuLuan_Lop10
\begin{bt}%[0D2H1-3]%[ma:a5d9014e]
	Trong $1$ lạng ($100$ gam) thịt bò chứa khoảng $26$ gam protein, $1$ lạng cá rô phi chứa khoảng $20$ gam protein. Trung bình trong một ngày, một người  phụ nữ cần tối thiểu $46$ gam protein. Gọi $x$, $y$ lần lượt là số lạng thịt bò và số lạng cá rô phi mà một người phụ nữ nên ăn trong một ngày. Viết bất phương trình bậc nhất hai ẩn $x$, $y$ để biểu diễn lượng protein cần thiết cho một người phụ nữ trong một ngày và chỉ ra ba nghiệm của bất phương trình đó.
	\loigiai{\textbf{Bước 1:} Biểu diễn lượng protein  có trong $x$ lạng thịt bò và $y$ lạng cá rô phi.
		Lượng proteintrong $x$ lạng thịt bò là $26x$  (gam).
		Lượng proteintrong $y$ lạng cá rô phi  là $20y$ (gam).
		Lượng proteintrong $x$ lạng thịt bò và $y$  lạng cá rô phi là $26x+20y$ (gam).\\
		\textbf{Bước 2:} Biểu diễn bất  phương trình.
		Vì  lượng proteintối thiểu là $46$ gam nên ta có bất phương trình $26x+20y \ge 46$.\\
		 \textbf{Bước 3:} Tìm nghiệm của bất phương trình.
		Thay $x=1$, $y=1$ vào bất phương trình ta được $26 \cdot 1+ 20 \cdot 1 \ge 46$. 
		Thay $x=2$,  $y=1$ vào bất phương trình ta được $26 \cdot 2+ 20 \cdot 1 \ge 46$.
		Thay $x=1$, $y=2$ vào bất phương trình ta được $26 \cdot 1+ 20 \cdot 2  \ge 46$.
		Vậy $(1;1)$, $(2;1)$, $(1;2)$ là các nghiệm  cần tìm.
	 }
\end{bt}

% nguồn: MaHoaNH_TuLuan_Lop10
\begin{bt}%[0D2H1-3]%[ma:0506d511]
	Hà, Châu, Liên và Ngân cùng đi mua trà sữa. Cả bốn bạn có tất cả $185$ nghìn đồng. Bốn bạn mua $4$ cốc trà sữa với giá tiền $35$ nghìn đồng một cốc. Các bạn gọi thêm trân  châu cho vào trà sữa. Một phần trân châu đen có giá $5$ nghìn đồng, một phần trân châu trắng có giá $10$ nghìn đồng. Gọi $x$, $y$ lần lượt là số phần trân châu đen, trân châu trắng mà bốn bạn định mua thêm.
	\begin{enumerate}
		\item Viết bất phương trình bậc nhất hai ẩn $x$, $y$ để  thể hiện số tiền các bạn có đủ khả năng chi trả cho phần trân châu đen, trắng.
		 \item Chỉ ra một nghiệm nguyên của bất phương trình đó.
	\end{enumerate}
	\loigiai{ 
		\begin{enumerate}
			 \item
			Tổng số tiền bốn bạn  phải trả để mua trân châu đen là $5x$ (nghìn đồng).\\
			Tổng số tiền bốn bạn phải trả để mua trân châu trắng  là $10y$ (nghìn đồng).\\
			Vì cả bốn bạn có tất cả 185 nghìn đồng mà các bạn đã dùng $35 \cdot 4 = 140$ nghìn đồng để mua trà sữa nên các bạn chỉ còn $45$  nghìn đồng.\\
			Do đó tổng  số tiền mua trân châu không thể vượt quá $45$ nghìn đồng nên ta có bất phương trình
			$$5x+10y \le 45 \;\, \text{hay}  \;\, x+2y \le 9.$$
			\item Vì  $4+ 2 \cdot 2 \le 9$ nên $(4;2)$ là một nghiệm nguyên của bất phương trình.
	\end{enumerate}}
\end{bt}

% nguồn: MaHoaNH_TuLuan_Lop10
\begin{bt}%[0D2H1-3]%[ma:862300b3]
	Để chào mừng năm học mới, mẹ An cho An $30$ đồng để mua thêm một số đồ dùng  học tập. Biết $1$ quyển sách nâng cao có giá $5$ đồng, $1$ quyển vở có giá $3$ đồng, $1$ cái bút có giá $2$ đồng. Gọi số lượng sách, vở và bút mà An mua lần lượt là $x$, $y$ và $z$.
	\begin{enumerate}
		\item Lập bất phương trình theo $x$ $y$, $z$ để diễn tả giới hạn về số lượng sách, vở và bút  mà An có thể mua được trong các trường hợp sau.
		 \textbf{Trường hợp 1:} An chỉ mua sách và vở.
		\textbf{Trường hợp 2:} An chỉ mua bút và vở. 
		\item Dùng bất  phương trình ở câu a) để trả lời hai câu hỏi sau
		\textbf{Trường hợp 1:} Nếu An mua $3$ quyển sách và $5$ quyển  vở thì có phù hợp không?
		\textbf{Trường hợp 2:} Nếu An mua $5$ quyển vở và  $8$ cái bút thì có phù hợp không?
	\end{enumerate}
	\loigiai{\begin{enumerate}
		\item Bất phương trình theo $x$, $y$, $z$  để diễn tả giới hạn về số lượng sách, vở và bút mà An có thể mua được trong từng trường hợp lần lượt là $5x+3y \le 30$ và $3y+2z \le 30$.
		\item \textbf{Trường hợp 1:} Số tiền nếu An mua $3$ quyển sách và $5$  quyển vở là $5 \cdot 3+ 3 \cdot 5=30 \le 30$ (thỏa mãn).\\
		Vậy trong trường hợp 1 thì An mua số lượng sách và vở  phù hợp với số tiền An có.
		\textbf{Trường hợp 2:} Số tiền nếu An mua $5$ quyển vở  và $8$ cái bút là $5 \cdot 3+ 2 \cdot 8 =31 \le 30$ (vô lý).
		 Vậy trong trường hợp 2 thì An mua số lượng sách và vở không phù hợp với số tiền An có.
\end{enumerate}}
\end{bt}

% nguồn: MaHoaNH_TuLuan_Lop10
\begin{bt}%[0D2H1-3]%[ma:926885f4]
	Cho biết mỗi $100$g thịt bò chứa $250$ calo, một quả trứng nặng $44$g chứa $70$ calo.  Giả sử có một người mỗi buổi sáng cần không quá $600$ calo. Gọi số gam thịt bò và số quả trứng mà người đó ăn trong một buổi sáng lần lượt là $x$ và $y$.
	\begin{enumerate}
		\item  Lập bất phương trình theo $x$, $y$ diễn tả giới hạn về lượng calo trong khẩu phần ăn buổi sáng của người đó.
		\item Dùng bất phương trình ở câu a) để trả lời  hai câu hỏi sau
	\end{enumerate}
	\begin{itemize}
		\item \textbf{Trường hợp 1:}  Nếu người đó ăn $200$g thịt bò và $2$ quả trứng (mỗi quả nặng $44$g) trong buổi sáng thì có phù hợp không?
		\item\textbf{ Trường hợp 2:} Nếu người đó ăn $150$g thịt bò và $3$ quả trứng (mỗi quả nặng $44$g) trong  buổi sáng thì có phù hợp không?
	\end{itemize}
	\loigiai{\begin{enumerate}
		\item Bất phương trình theo $x$, $y$ diễn tả giới hạn  về lượng calo trong khẩu phần ăn buổi sáng của người đó là
		 $$\dfrac{x}{100} \cdot 250 +70y \le 600 \Leftrightarrow 2{,}5x+70y \le 600.$$
		\item \textbf{Trường hợp  1:} Lượng calo hấp thụ được là $2{,}5 \cdot 200+ 70 \cdot 2 =640 \le 600$ (vô lý).\\
		Vậy trong trường hợp 1 không phù hợp với yêu cầu đề  ra.
		\textbf{Trường hợp 2:} Lượng calo hấp thụ được là $2{,}5 \cdot 150+ 70 \cdot 3 =585 \le 600$ (thỏa  mãn).\\
		Vậy trong trường hợp 2 phù hợp với yêu cầu  đề ra.
\end{enumerate}}
\end{bt}

% nguồn: MaHoaNH_TuLuan_Lop10
\begin{bt}%[0D2H1-3]%[ma:6ea1df3a]
	Một công ty sản xuất hai loại sản phẩm gồm sản phẩm $X$ và sản phẩm $Y$. Chi phí sản xuất một sản phẩm $X$ là $2$ triệu đồng, sản xuất một sản phẩm $Y$  là $3$ triệu đồng. Chi phí thuê mặt bằng kho chứa là $100$ triệu đồng.
	\begin{enumerate}
		\item Tính tổng chi  phí công ty sản xuất $200$ sản phẩm loại $X$ và $150$ sản phẩm loại $Y$.
		\item Gọi $x$ và $y$ lần lượt là số sản phẩm $X$ và sản phẩm $Y$ công ty sản xuất.  Biết vốn sản xuất của công ty là $1$ tỉ đồng. Viết bất phương trình biểu thị mối liên hệ giữa $x$ và $y$.
	\end{enumerate}
	 \loigiai{
		\begin{enumerate}
			\item Tổng chi phí công ty  sản xuất $200$ sản phẩm loại $X$ và $150$ sản phẩm loại $Y$ là
			$$200\cdot2+150\cdot3+100=950 \, \text{triệu đồng.}$$ 
			\item  Chi phí sản xuất $x$ sản phẩm $X$ là $2x$ triệu đồng.\\
			Chi phí sản xuất $y$  sản phẩm $Y$ là $3x$ triệu đồng.\\
			Chi  phí thuê mặt bằng kho chứa là $100$ triệu đồng.\\
			Tổng chi phí là $2x+3y+100$ triệu  đồng.\\
			Vì vốn sản xuất của công  ty là $1$ tỉ đồng $=$ $1\,000$ triệu đồng.\\
			Do đó bất phương trình biểu  thị mối liên hệ giữa $x$ và $y$ là
			$$2x+ 3y +100\le 1\,000\Leftrightarrow 2x+3y\le  900.$$
			Vậy bất  phương trình cần tìm là $2x+3y\le 900$.
		\end{enumerate}
	 }
\end{bt}

% nguồn: MaHoaNH_TuLuan_Lop10
\begin{bt}%[0D2H1-3]%[ma:1fce2252]
	Ông An muốn thuê một chiếc ô tô (có lái xe) trong một tuần. Giá thuê xe được cho như bảng sau: 
	\begin{center}
		\begin{tabular}{|c|c|c|}
			\hline 
			&  Phí cố định (nghìn đồng/ngày) &
			\begin{tabular}{c}
				 Phí tính theo quãng đường di\\
				chuyển  (nghìn đồng/km)\\
			\end{tabular} \\
			\hline 
			Từ thứ $2$ đến thứ $6$  & $900$ & $8$\\
			\hline 
			Thứ $7$ và Chủ nhật& $1\,500$ & $10$\\ 
			 \hline
		\end{tabular}
	\end{center}
	\begin{enumerate}
		\item Biết các ngày từ thứ Hai đến thứ Sáu mỗi ngày ông An di chuyển $12$ km và  ngày Chủ nhật ông An di chuyển $60$ km. Tính số tiền ông An phải trả.
		\item Gọi $x$ và $y$ lần lượt là số kilômét ông An đi trong các ngày từ thứ Hai đến thứ Sáu và trong hai ngày cuối tuần. Viết bất phương trình  biểu thị mối liên hệ giữa $x$ và $y$ sao cho tổng số tiền ông An phải trả không quá $14$ triệu đồng.
	\end{enumerate}
	\loigiai{ 
		\begin{enumerate}
			\item Số tiền ông An đi $12$ km trong các ngày từ thứ Hai đến thứ Sáu là  $900\cdot5 +8\cdot12=4\,596$ (nghìn đồng).\\
			 Số tiền ông An đi $60$ km trong Chủ nhật là $1\,500\cdot1 + 10\cdot60=2\,100$ (nghìn đồng).\\
			 Tổng số tiền là $4\,596 + 2\,100=6\,696$ ngàn đồng.
			\item Ta có $14$ triệu đồng $=$ $1\,400$ nghìn đồng.\\ 
			Số tiền ông An đi $x$ km trong các ngày từ thứ Hai  đến thứ Sáu là $900\cdot5 +8x$ (nghìn đồng).\\
			Số tiền ông An đi $y$ km trong $2$ ngày thứ Bảy và Chủ nhật là $1\,500\cdot2 + 10y$  (nghìn đồng).\\
			Do đó số tiền ông An đi trong một tuần là $8x + 10 y  +7\,500$ (nghìn đồng).\\
			Vì số tiền không quá $14$  triệu đồng nên ta có
			\[8x + 10 y +7\,500\le 1\,400 \Leftrightarrow 4x+  5y \le 3\,250.\]
			Vậy bất phương trình cần tìm là $4x+ 5y  \le 3\,250$.
		\end{enumerate}
	 }
\end{bt}

% nguồn: MaHoaNH_TuLuan_Lop10
\begin{bt}%[0D2H1-3]%[ma:694a412b]
	Một cửa hàng bán lẻ bán hai loại hạt cà phê. Loại thứ nhất giá $140$ nghìn đồng/kg và loại thứ hai giá $180$ nghìn đồng/kg. Cửa hàng trộn $x$ kg loại thứ nhất  và $y$ kg loại thứ hai sao cho hạt cà phê đã trộn có giá không quá $170$ nghìn đồng/kg.
	\begin{enumerate}
		\item Viết bất phương trình bậc nhất hai ẩn $x$,  $y$ thoả mãn điều kiện đề bài.
		\item Biểu diển  miền nghiệm của bất phương trình tìm được ở câu a trên mặt phẳng toạ độ.
	\end{enumerate}
	 \loigiai{
		\begin{enumerate}
			 \item
			 Giá bán $x$ kg cà phê loại thứ nhất là $140x$ (nghìn đồng).\\
			Giá bán $y$  kg cà phê loại thứ hai là $180x$ (nghìn đồng).\\
			 Giá bán khi trộn $x$ kg loại thứ nhất và $y$ kg loại thứ hai là $140x+180y$ (nghìn đồng).\\
			Theo đề  bài ta có $140x+180y\le 170(x+y)\Leftrightarrow 3x-y\ge 0$.\\
			 Vậy bất phương trình bậc nhất hai ẩn $x$, $y$ thoả mãn điều kiện đề bài là $3x-y\ge 0$.
			\item Biểu diển miền nghiệm của bất phương trình $3x-y\ge  0$.
			\begin{itemize}
				\item Vẽ đường thẳng $\Delta\colon  3x-y= 0$ đi qua hai điểm $\left(0;0\right)$ và $\left(1;3\right)$.
				\item Xét điểm $M(-1;1)\notin \Delta$. Ta có  $3\cdot(-1)-1=-4<0$.
				\item Do đó, miền nghiệm của bất phương trình $3x-y\ge 0$ là nửa mặt phẳng bờ $\Delta$ không chứa điểm $M$ kể cả đường  thẳng $\Delta$.
			\end{itemize}
			\begin{center}
				\begin{tikzpicture}[scale=.9,>=stealth,font=\footnotesize,line join=round,line cap=round]
						\def\a{3}  \def\b{0}
						 \def\xmin{-3.5} \def\xmax{2}
						\def\ymin{-2}  \def\ymax{4}
						\draw[->]  (\xmin,0)--(\xmax+.2,0) node [below]{$x$};
						\draw[->] (0,\ymin)--(0,\ymax+.2) node [left]{$y$}; 
						\fill (0,0)  circle (1pt) node[below left]{$O$};
						\clip (\xmin,\ymin)  rectangle (\xmax,\ymax);
						 \fill[pattern=north east lines](\xmin,\ymax)--(\xmin,\ymin)--plot (\x,{\a*(\x)+\b})--cycle;
						\draw[smooth,samples=300]  plot (\x,{\a*(\x)+\b});
						 \node[rotate=72,right] at (.5,.5) {$\Delta\colon 3x-y =0$};
						\draw[fill=black] (-1,1) circle  (1pt) node [above left,fill=white,rectangle]{$M$};
						\draw[dashed] (-1,0) node[below]{$-1$}|-(0,1)  node[right]{$1$};
					\end{tikzpicture}
			\end{center}
		\end{enumerate}
	 }
\end{bt}

% nguồn: MaHoaNH_TuLuan_Lop10
\begin{bt}%[0D2H1-3]%[ma:6436a4f1]
	Ông An muốn thuê một chiếc ô tô (có lái xe) trong một tuần. Giá thuê xe được cho  như bảng sau:
	\begin{center}
		\begin{tabular}{|c|c|c|}
			 \hline
			& Phí cố định (nghìn  đồng/ngày) &
			\begin{tabular}{c}
				Phí tính theo quãng  đường di\\
				 chuyển (nghìn đồng/km)\\
			\end{tabular} \\
			\hline 
			Từ  thứ $2$ đến thứ $6$ & $900$ & $8$\\
			\hline 
			Thứ $7$ và  Chủ nhật& $1\,500$ & $10$\\
			\hline 
		\end{tabular}
	\end{center}
	\begin{enumerate}
		\item Gọi $x$ và $y$ lần lượt là số kilômét ông An đi trong các ngày từ thứ Hai đến thứ Sáu và trong hai ngày cuối  tuần. Viết bất phương trình biểu thị mối liên hệ giữa $x$ và $y$ sao cho tổng số tiền ông An phải trả không quá $14$ triệu đồng.
		\item Biểu diễn miền nghiệm của bất phương trình  trên mặt phẳng toạ độ.
	\end{enumerate}
	 \loigiai{
		\begin{enumerate}
			\item Ta có $14$ triệu đồng $=$ $1\,400$  nghìn đồng.\\
			Số tiền ông An đi  $x$ km trong các ngày từ thứ Hai đến thứ Sáu là $900\cdot5 +8x$ (nghìn đồng).\\
			Số tiền ông An đi $y$ km trong $2$ ngày thứ Bảy và Chủ nhật là $1\,500\cdot2 + 10y$  (nghìn đồng).\\
			Do đó số tiền  ông An đi trong một tuần là $8x + 10 y +7\,500$ (nghìn đồng).\\
			Vì số  tiền không quá $14$ triệu đồng nên ta có
			\[8x + 10 y +7\,500\le 1\,400 \Leftrightarrow 4x+ 5y \le 3\,250.\] 
			Vậy bất phương trình cần tìm là $4x+  5y \le 3\,250$.
			\item Biểu  diễn miền nghiệm của bất phương trình $4x+ 5y \le 3\,250$.
			\begin{itemize}
				\item Vẽ đường thẳng $\Delta\colon 4x+ 5y = 3\,250$ đi qua  hai điểm $\left(812{,}5;0\right)$ và $\left(0;650\right)$.
				\item Xét điểm $O(0; 0)\notin \Delta$. Ta  có $4\cdot0+5\cdot0=0<3250$.
				\item Do đó, miền nghiệm của bất phương trình $4x+ 5y \le 3\,250$ là nửa mặt phẳng bờ $\Delta$ chứa điểm $O$ kể cả đường thẳng  $\Delta$.
			\end{itemize}
			\begin{center}
				\begin{tikzpicture}[scale=.9,>=stealth,font=\footnotesize,line join=round,line cap=round]
						 \def\a{-0.7} \def\b{2.1}
						\def\xmin{-3.5} \def\xmax{6} 
						\def\ymin{-1} \def\ymax{5} 
						\draw[->]  (\xmin,0)--(\xmax+.2,0) node [below]{$x$};
						\draw[->] (0,\ymin)--(0,\ymax+.2)  node [left]{$y$};
						\fill (0,0) circle (1pt)  node[below left]{$O$};
						 \clip (\xmin,\ymin) rectangle (\xmax,\ymax);
						\fill[pattern=north east lines]plot  (\x,{\a*(\x)+\b}) --(\xmax,\ymin) --(\xmax,\ymax)--(\xmin,\ymax) --cycle;
						\draw[smooth,samples=300] plot (\x,{\a*(\x)+\b}); 
						\node[rotate=-35,right] at  (-3.5,4) {$\Delta\colon 4x+ 5y = 3\,250$};
						\draw[fill=black] (0,2.1)  circle (1pt) node [yshift=-10pt,xshift=-10pt]{$650$};
						\draw[fill=black] (3,0) circle (1pt) node[yshift=-12pt,xshift=-5pt]{$812{,}5$}; 
					\end{tikzpicture}
			\end{center}
		\end{enumerate}
	 }
\end{bt}
